\setchapterpreamble[u]{\margintoc}
\chapter{T-SQL: Controle de Dados (DCL)}
\labch{cet615-dcl}

Este capítulo cobre a última apresentação da
disciplina~\cite{BaiocoCET615}: controle de acesso, cópia de
segurança e agendamento de tarefas. A aula tratou o tema como uma visão
geral: quem já trabalha como DBA conhece, e quem não conhece fica sabendo
o que existe para pesquisar depois. Não houve demonstração ao vivo; as
telas do SSMS são as dos slides.

\section{Por que usar o controle do SGBD}

Muitos sistemas fazem o controle de acesso só na aplicação: uma tabela
de usuários e senhas dentro do banco e uma conexão que pode tudo. A
professora chamou isso de ``bem temeroso'' por causa da \textbf{injeção
de SQL}\index{injeção de SQL} (\emph{SQL injection}): com trechos de
comando digitados no campo de login, forçando erros não tratados,
alguém descobre um acesso válido, e por trás está uma conexão sem
restrição. A recomendação é usar o controle de acesso do próprio SGBD,
que nos fabricantes estabelecidos é bem elaborado e serve a aplicações em
qualquer linguagem.

\section{Três níveis de acesso}

Para ler um dado, um usuário precisa passar por três portas:

\begin{enumerate}
	\item \textbf{Servidor:} um \emph{login} permite conectar à instância.
	\item \textbf{Banco:} um \emph{user} (usuário do banco), ligado ao
	login, permite entrar em um banco específico.
	\item \textbf{Objetos:} \emph{permissões} dizem o que o usuário pode
	fazer com cada tabela, visão ou procedimento.
\end{enumerate}

Um login sem usuário no banco conecta ao servidor mas não abre o banco;
um usuário sem permissão na tabela abre o banco mas não lê a tabela. O
\texttt{sa} (\emph{system administrator}) é um login com todas as
permissões; um login comum pode servir só para conectar.

\begin{figure}[h]
\centering
\begin{tikzpicture}
	\node[caixa, fill=white, minimum width=20mm] (p) at (0,0) {pessoa ou\\aplicação};
	\node[draw=bdroxo, fill=bdroxoclaro, rounded corners, minimum width=74mm, minimum height=36mm] (srv) at (5.6,0) {};
	\node[rotulo, anchor=north west] at (srv.north west) {servidor: login};
	\node[draw=bdazul, fill=bdazulclaro, rounded corners, minimum width=54mm, minimum height=24mm] (bd) at (6.3,-0.4) {};
	\node[rotulo, anchor=north west] at (bd.north west) {banco: usuário};
	\node[caixa, fill=white, minimum width=36mm] (obj) at (6.6,-0.8) {objetos: permissões\\(select, execute ...)};
	\draw[seta] (p) -- node[rotulo, above] {conecta} (srv.west);
\end{tikzpicture}
\caption{Os três níveis de acesso, de fora para dentro.}
\end{figure}

\section{Logins e usuários}

\begin{lstlisting}
use master
go
create login leitor_app
with password = 'Senha#Forte2026', default_database = biblioteca
go
alter login leitor_app with password = 'Outra#Senha2026'
drop login leitor_app

use biblioteca
go
create user app_biblioteca for login leitor_app
drop user app_biblioteca
\end{lstlisting}

O \texttt{create login} roda no \texttt{master}, onde fica o catálogo
de logins, e o \texttt{create user}, no banco de destino. O usuário não
tem senha: a senha é do login, e o usuário só o mapeia para dentro do
banco. Sem \texttt{default\_database}, o banco padrão do login é o
\texttt{master}, o que a professora considera ruim. Os procedimentos antigos (\texttt{sp\_addlogin},
\texttt{sp\_adduser}, \texttt{sp\_droplogin} e outros) foram substituídos
por esses comandos a partir do SQL Server 2005. Tudo isso também pode ser
feito pelo SSMS: a pasta \emph{Security} do servidor tem logins e papéis
de servidor, e a pasta \emph{Security} de cada banco tem usuários,
papéis e esquemas.

\section{Papéis (\emph{roles})}

Controle de acesso é projeto, ``não é assim do nada''. Um
\textbf{papel}\index{role} é um grupo de usuários com outro nome: agrupa
quem precisa das mesmas permissões.
Quem entra no papel herda as permissões dele, o que simplifica a
administração. Há dois tipos:

\begin{description}
	\item[Papéis de servidor:] agrupam logins e dão permissões
	administrativas. O \texttt{sysadmin} pode tudo na instância; outros
	cobrem partes do trabalho, como \texttt{securityadmin} (logins),
	\texttt{serveradmin} (opções do servidor) e \texttt{dbcreator}
	(criação de bancos). A professora sugere que o DBA os use para delegar
	tarefas, ou para separar as próprias, em vez de fazer tudo como
	\texttt{sa}. Desde o SQL Server 2012 é possível criar papéis de
	servidor próprios (\texttt{create server role}); ela prefere combinar
	os papéis prontos.
	\item[Papéis de banco:] agrupam usuários de um banco. Há papéis
	prontos, como \texttt{db\_owner} (controle total do banco),
	\texttt{db\_datareader}, \texttt{db\_datawriter} e
	\texttt{db\_ddladmin}, que permite criar, alterar e excluir objetos e
	não deve ir para o usuário final. É possível criar novos com
	\texttt{create role}. Todo usuário pertence automaticamente ao papel
	\texttt{public}.
\end{description}

\begin{lstlisting}
alter server role sysadmin add member leitor_app   -- login
alter server role sysadmin drop member leitor_app

use biblioteca
create role atendentes
alter role atendentes add member app_biblioteca    -- usuário
alter role atendentes drop member app_biblioteca
drop role atendentes
\end{lstlisting}

\subsection*{Complemento: o que o papel \texttt{public} permite}

O slide diz que o \texttt{public} tem permissões limitadas a leitura de
dados, e em aula a professora alertou que, por isso, um usuário recém
criado já poderia fazer \texttt{select}. A documentação diz que o
\texttt{public} não tem permissões próprias sobre os dados: ele lê
metadados do sistema\sidecite{MSDatabaseRoles}, e o usuário herda dele só o que alguém conceder ao
\texttt{public}. No teste, um usuário novo, só no \texttt{public},
recebeu \emph{The SELECT permission was denied} ao consultar uma tabela.
O alerta continua útil: o que for concedido ao \texttt{public} vale para
todos.

\section{Permissões}

Permissões definem quem pode fazer o quê, em três níveis:

\begin{description}
	\item[Servidor:] vêm principalmente dos papéis de servidor; desde o
	SQL Server 2012 também podem ser dadas diretamente a um login.
	\item[Banco:] criar objetos, fazer backup do banco ou do log. Vêm dos
	papéis de banco ou de \texttt{grant} direto.
	\item[Objeto:] consultar, inserir, alterar e excluir dados de uma
	tabela, executar um procedimento.
\end{description}

\begin{lstlisting}
-- banco
grant create table, create view to app_biblioteca
revoke create table, create view from app_biblioteca

-- objeto
grant select, insert on emprestimo to atendentes
grant execute on ins_livro_exemplar to atendentes
revoke insert on emprestimo from atendentes
deny delete on emprestimo to atendentes
\end{lstlisting}

Na prática, as permissões de banco são de DDL e as de objeto, de DML.
As de objeto chegam até a coluna, mas fazer isso com \texttt{grant} um a
um é trabalhoso, e é aí que os papéis ajudam.

\begin{description}
	\item[\texttt{grant}] concede a permissão.
	\item[\texttt{revoke}] retira uma permissão concedida \emph{ou
	negada} antes; o usuário volta ao estado neutro e pode ainda receber
	a permissão por um papel.
	\item[\texttt{deny}] nega explicitamente. A negação tem precedência:
	quem tem \texttt{deny} não recebe a permissão nem por papel.
\end{description}

A professora deu dois usos para o \texttt{deny}: negar temporariamente
sem retirar a permissão, durante uma investigação, por exemplo; e criar
um papel só de negações para os dados sensíveis, no qual entra todo
usuário novo. Combinar um papel que concede com outro que nega também
evita criar um papel para cada combinação, como um papel ``só
\texttt{insert}''.

O slide escreve \texttt{revoke ... to}; a documentação aceita
\texttt{to} e \texttt{from} nesse comando\sidecite{MSRevoke}, e
\texttt{from} lê melhor.

Este é o mecanismo por trás da ``segurança por funcionalidade'' dos
capítulos de visões e procedimentos: em vez de dar \texttt{select} nas
tabelas, dá-se \texttt{select} na visão ou \texttt{execute} no
procedimento, e o usuário nunca toca as tabelas. A professora insistiu
nisso: uma visão que já restringe linhas e colunas, ou um procedimento
que faz uma transação em várias tabelas, dá ao usuário exatamente a
funcionalidade de que ele precisa. Ela acha esses recursos pouco usados,
enquanto se tenta resolver tudo na interface, com brechas.

\subsection*{Complemento: \texttt{grant all}}

Em aula, \texttt{grant all} no banco foi comparado a pôr o usuário no
\texttt{db\_owner}. A documentação diz outra coisa: \texttt{all} está
obsoleto, existe só por compatibilidade e não concede todas as
permissões. No banco, equivale a \texttt{backup database},
\texttt{backup log}, \texttt{create database}, \texttt{create default},
\texttt{create function}, \texttt{create procedure}, \texttt{create
rule}, \texttt{create table} e \texttt{create view}\sidecite{MSGrant}.
O \texttt{db\_owner} tem o controle total do banco.

\subsection*{Complemento: exceções à precedência do \texttt{deny}}

A documentação registra uma exceção mantida por compatibilidade: um
\texttt{deny} na tabela inteira não prevalece sobre um \texttt{grant} em
uma coluna específica dela\sidecite{MSRevoke}. Também registra que a
forma \texttt{revoke all} está obsoleta; o certo é revogar permissões
específicas.

\subsection*{Complemento: propagação de permissões}

\textcite[seção 24.2]{Elmasri2011} trata o controle de acesso
discricionário do padrão SQL. Quem recebe uma permissão \texttt{with
grant option} pode repassá-la a outros. Ao revogar, o SGBD precisa
retirar também as permissões que foram repassadas a partir dela. No SQL
Server, isso exige \texttt{revoke grant option for ... cascade}; sem o
\texttt{cascade}, o \texttt{revoke} falha quando a permissão foi
concedida com \texttt{grant option}\sidecite{MSRevoke}. O mesmo capítulo
apresenta o controle de acesso baseado em papéis (seção 24.3), que é o
modelo dos \emph{roles} do SQL Server.

\section{Esquemas}

Um \textbf{esquema}\index{esquema (schema)} (\emph{schema}) agrupa
objetos de um banco que têm o mesmo dono, um usuário ou um papel. Os
donos passam a ser donos de esquemas, e não de cada objeto. Permissões
podem ser dadas no esquema inteiro. Dentro de um esquema, nomes de
objetos não se repetem. O nome completo de um objeto é
\texttt{servidor.banco.esquema.objeto}. Trocar o dono de um esquema não
muda esse nome.

O uso que a professora destacou: quando se decide dar acesso direto às
tabelas, sem visão nem procedimento, concede-se uma vez no esquema que
agrupa as tabelas, em vez de repetir o \texttt{grant} em cada uma. Ela
vê o melhor uso num banco central de empresa com departamentos bem
separados.

\begin{lstlisting}
grant select on schema :: dbo to atendentes
\end{lstlisting}

\subsection*{Complemento: comandos de esquema}

Em aula, ficou a dúvida se há comandos para criar esquemas e mover
objetos; há. O teste abaixo rodou no SQL Server 2022 e terminou com
\texttt{dbo.livro}, \texttt{acervo.livro} e \texttt{acervo.exemplar}
no mesmo banco:

\begin{lstlisting}
-- create schema precisa ser o primeiro comando do lote
create schema acervo
go
-- mesmo nome de dbo.livro, em outro esquema
create table acervo.livro (codlivro int)
-- move a tabela para o esquema acervo
alter schema acervo transfer dbo.exemplar
-- troca o dono do esquema
alter authorization on schema :: acervo to app_biblioteca
\end{lstlisting}

O teste confirma a regra do slide: o nome só precisa ser único dentro
do esquema, e dois esquemas do mesmo banco podem ter tabelas com o mesmo
nome. O esquema pertence a um banco e não agrupa objetos de bancos
diferentes.

\section{Cópia de segurança e recuperação}

Backup não é copiar os arquivos \texttt{.mdf}, \texttt{.ldf} e
\texttt{.ndf} pelo sistema operacional. Ele é feito de dentro do
servidor, porque a estrutura está no catálogo, que também precisa de
cópia. Como na modelagem, a estratégia depende dos requisitos e do valor
da informação: o banco precisa ficar no ar
24 horas por dia? Quanto tempo de parada se aceita depois de uma falha?
As possibilidades citadas na aula:

\begin{description}
	\item[Espelhamento (\emph{mirroring}):] os dados ficam duplicados em
	dois servidores, de preferência em lugares diferentes; se um falha, o
	outro assume em segundos ou minutos. Para a professora, é ``o melhor
	dos mundos'' e a única saída para bancos que mudam o tempo todo, como
	os de um banco comercial; uma loja pode viver com instantâneos ou
	backup completo.
	\item[Instantâneos (\emph{snapshots}):] cópias do banco em intervalos
	curtos.
	\item[Backup e restore:] cópias para restauração posterior.
\end{description}

\subsection{Tipos de backup}

\begin{description}
	\item[Completo (\emph{full}):] o banco inteiro. É o mais lento.
	\item[Diferencial:] só as páginas alteradas desde o último completo.
	\item[Do log de transações:] só o log.
\end{description}

Duas estratégias básicas:

\begin{itemize}
	\item \textbf{Só o banco}, pelo menos diariamente. O que mudou desde o
	último backup se perde numa falha.
	\item \textbf{Banco e log.} O completo é menos frequente (por exemplo,
	semanal), e o log é copiado várias vezes por dia. Na recuperação,
	restaura-se o último completo e depois todos os backups de log, em
	ordem cronológica. O diferencial encurta essa sequência.
\end{itemize}

O exemplo da aula para o backup só do banco: um comércio que funciona
oito horas por dia pode aceitar ficar o dia descoberto e fazer o backup
completo no fim do expediente. O objetivo de qualquer estratégia é
cobrir o maior período possível de alterações.

O \texttt{master} deve ter backup regular e também logo depois de criar,
alterar ou excluir bancos, logins, usuários, papéis ou permissões.
``Regular'' quer dizer que, sem essas mudanças, ele fica parado; depois
delas, o backup é imediato, ou se perde o catálogo. A professora
recomenda também um backup completo do banco de usuário depois de
mudanças de estrutura.

\begin{lstlisting}
backup database biblioteca to disk = 'C:\backup\biblioteca.bak'
backup log biblioteca to disk = 'C:\backup\biblioteca_log.trn'

restore database biblioteca from disk = 'C:\backup\biblioteca.bak'
   with norecovery          -- ainda virão logs
restore log biblioteca from disk = 'C:\backup\biblioteca_log.trn'
   with recovery            -- último passo: o banco volta ao ar
\end{lstlisting}

No SSMS: botão direito no banco $\rightarrow$ \emph{Tasks} $\rightarrow$
\emph{Back Up} ou \emph{Restore}. A aula cita duas novidades do SQL Server
2005: a restauração parcial, que deixa o banco disponível assim que o
grupo de arquivos primário é restaurado, e o backup do log em paralelo
com o backup completo.

\subsection*{Complemento: backup com o banco no ar}

Em aula, foi dito que backup e restauração podem exigir tirar o banco do
ar. No SQL Server, os backups completo, diferencial e de log rodam com o
banco on-line e em uso; o custo é de leitura em disco e desempenho. No
teste, um \texttt{backup database} rodou com uma conexão usando o
banco. Quem costuma deixar o banco indisponível é a restauração, e é por
isso que a restauração parcial foi tratada como novidade.

Também apareceu a ideia de que o \emph{checkpoint} é o momento de
esvaziar o log. São coisas diferentes: o checkpoint grava em disco as
páginas alteradas que estão na memória. No modelo de recuperação
completo, o espaço do log só é liberado depois do backup do log; no
simples, depois do checkpoint\sidecite{MSTransactionLog}.

\subsection*{Complemento: modelos de recuperação}

O que pode ser copiado depende do \textbf{modelo de recuperação} do
banco\sidecite{MSRecoveryModels}:

\begin{description}
	\item[Simples:] não permite backup de log. O SQL Server reaproveita o
	espaço do log sozinho. Numa falha, perde-se o que mudou desde o
	último backup.
	\item[Completo (\emph{full}):] exige backup de log, e o log cresce até
	que ele seja feito. Permite restaurar até um ponto exato no tempo,
	por exemplo antes de um erro de um usuário.
	\item[\emph{Bulk-logged}:] variação do completo que registra pouco as
	cargas em massa; não permite restaurar até um ponto no tempo.
\end{description}

As edições Enterprise e Standard usam o modelo completo por padrão; a
Express usa o simples. A estratégia ``banco e log'' da aula só funciona
no modelo completo ou no \emph{bulk-logged}. Isso explica também o
alerta do \refch{cet615-sqlserver} sobre logs que crescem demais: no
modelo completo, sem backup de log, o log não é truncado.

\subsection*{Complemento: o espelhamento está obsoleto}

A documentação marca o espelhamento de banco para remoção em versão
futura e recomenda os grupos de disponibilidade \emph{Always On} para
alta disponibilidade\sidecite{MSMirroring}. O espelhamento também só
funciona com o modelo de recuperação completo. Já o instantâneo de banco
do SQL Server é uma visão somente leitura, na mesma instância e
dependente do banco de origem: não substitui o backup.

\section{Agendamento de tarefas (\emph{jobs})}

Um \textbf{job}\index{job} é uma tarefa administrativa, com um ou mais
passos, executada automaticamente em datas e horários definidos: backups,
manutenção de índices e outras rotinas do DBA. Também pode ser criado
sem agendamento, para execução manual. Os jobs dependem do serviço
\emph{SQL Server Agent} e ficam guardados no banco \texttt{msdb}
(\refch{cet615-sqlserver}). No SSMS: \emph{SQL Server Agent}
$\rightarrow$ \emph{Jobs} $\rightarrow$ \emph{New Job}, com as abas
\emph{Steps} (passos) e \emph{Schedules} (agendamento). Os passos podem
usar os próprios comandos \texttt{backup} e \texttt{restore}. Mesmo sem
agendamento, o job pronto facilita a vida do DBA.

Com os jobs, a professora encerrou o conteúdo da disciplina. Segurança e
administração devem voltar com mais profundidade em outras disciplinas
do curso.
